\documentclass[11pt,a4paper]{article}
\usepackage{amsmath,amssymb,amsthm}
\usepackage{listings}
\usepackage{xcolor}
\usepackage{geometry}
\geometry{margin=1in}

\lstset{
  language=Lean,
  basicstyle=\ttfamily\small,
  keywordstyle=\color{blue},
  commentstyle=\color{green!50!black},
  stringstyle=\color{red},
  breaklines=true,
  frame=single,
  numbers=left
}

\title{\textbf{Poincaré Conjecture as a dm³ System (N=3 Exploration)}\\
A Lean-First Language Specification \\
(Topological Pillar of the Unified Framework)}
\author{Pablo Nogueira Grossi \\ G6 LLC / AXLE Project}
\date{v1.1 — April 2026}

\begin{document}

\maketitle

\begin{abstract}
N=3 exploration of the Poincaré Conjecture in dm³_poincare. Ricci flow is the generative step.
C = volume compression, K = scalar curvature, F = surgery, U = convergence to S³.
All operators and contact geometry are identical to the other five pillars.
\end{abstract}

\section{dm³ Grammar (Lean-first)}

\subsection{Syntax}
\begin{lstlisting}
inductive Dm3Op
| C | K | F | U
\end{lstlisting}

\subsection{N=3 Interpretation}
C = volume compression under Ricci flow,  
K = scalar curvature deformation,  
F = topological folding via surgery,  
U = unfolding to the 3-sphere.

\subsection{Composition}
\[
G = U \circ F \circ K \circ C
\]

\section{The dm³\_poincare Category (N=3)}

\subsection{Objects}
\begin{lstlisting}
structure Dm3PoincareObject :=
  (M            : ThreeManifold)
  (contactForm  : ThreeManifold → ThreeManifold)
  (simplyClosed : isSimplyConnectedClosed3 M)
  (dimension    : ℕ := 3)
\end{lstlisting}

\section{Poincaré as a dm³ Object (N=3)}

\subsection{State Space}
\begin{lstlisting}
def Poincare_state := ThreeManifold
\end{lstlisting}

\subsection{Object Definition}
\begin{lstlisting}
def Poincare_dm3 : Dm3PoincareObject := { ... }
\end{lstlisting}

\section{The Poincaré → dm³\_poincare Embedding (N=3)}

\subsection{Operator Decomposition}
\begin{lstlisting}
theorem Poincare_operatorDecomposition :
  ∀ M, Poincare_step M = (G C_Poincare K_Poincare F_Poincare U_Poincare) M := ...
\end{lstlisting}

\subsection{Contact Preservation}
\begin{lstlisting}
axiom Poincare_preserves_contact :
  ∀ M, Poincare_step (Poincare_contact M) = Poincare_contact (Poincare_step M)
\end{lstlisting}

\section{Remaining Analytic Axioms (N=3 Poincaré)}

\subsection{meanContraction\_poincare}
\begin{lstlisting}
axiom meanContraction_poincare :
  ∀ M, (∫ log (volume (Poincare_step M) / volume M) dμ) < 0
\end{lstlisting}

\subsection{lyapunovDescent\_poincare}
\begin{lstlisting}
axiom lyapunovDescent_poincare :
  ∀ M, scalarCurvature (Poincare_step M) < scalarCurvature M
\end{lstlisting}

\subsection{hasStructuredCycle\_poincare}
\begin{lstlisting}
axiom hasStructuredCycle_poincare :
  ∃ A, is_dm3_poincare_cycle A
\end{lstlisting}

\section{Coherence Bridge Synchronization}

\begin{lstlisting}[language=YAML]
- id: poincare_conjecture
  claim_level: formally_stated
  lean4_sorry_label: "ADMIT-D_poincare, ADMIT-E_poincare, ADMIT-F_poincare"
  lean4_sorry_count: 3
  axle_status: "Dm3Poincare.lean v1.1 — N=3 Exploration"
\end{lstlisting}

\section{Final Notes (N=3)}

In dimension N=3 the dm³ operators acquire concrete geometric meaning via Ricci flow:
C = volume decrease, K = scalar curvature evolution, F = surgery at singularities,
U = convergence to the 3-sphere. This makes the Poincaré Conjecture the natural topological
pillar of the unified dm³ framework. The three analytic admits remain the only gaps.

\end{document}
